% Header: General study / work / project template
%%%%%%%%%%%%%%%%%%%%%%%%%%%%%%%%%%%%%%%%%%%%%%%%%%%%%%%%%%%%%%%%%%%

\documentclass[11pt,a4paper]{scrartcl}

% Layout
\usepackage[margin=2.5cm]{geometry}
\usepackage[onehalfspacing]{setspace}

% General packages
\usepackage{graphicx}
\usepackage{enumitem}
\usepackage{tcolorbox}
\tcbuselibrary{breakable}
\usepackage[breaklinks=true,colorlinks=true,linkcolor=blue,urlcolor=blue,citecolor=blue]{hyperref}
\usepackage{amsmath,amsthm,amssymb,mathtools}
\usepackage{icomma}
\usepackage[T1]{fontenc}
\usepackage[utf8]{inputenc}

% Language
\usepackage[english,ngerman]{babel}

% Units / tables / floats / references
\usepackage[locale = UK,space-before-unit=true,per-mode = symbol]{siunitx}
\usepackage{booktabs,multirow}
\usepackage{placeins}
\usepackage[natbib,abbreviate=true,doi=false,style=numeric-comp,giveninits=true,sorting=none]{biblatex}
\usepackage{csquotes}

% Optional bibliography
% \addbibresource{references.bib}

% Graphics paths
\graphicspath{{Bilder/} {images/} {figures/} {chapters/}}

% Optional math shortcuts
\newcommand{\R}{\mathbb{R}}
\newcommand{\N}{\mathbb{N}}
\newcommand{\Q}{\mathbb{Q}}
\newcommand{\set}[1]{\left\{ #1 \right\}}
\newcommand{\abs}[1]{\left| #1 \right|}
\newcommand{\norm}[1]{\left\lVert #1 \right\rVert}
\newcommand{\ol}[1]{\overline{#1}}

% Optional units
\DeclareSIUnit{\dBm}{dBm}

% Box styles
\newtcolorbox{calculationbox}[1][]{
    title=Calculation / Rechnung,
    breakable,
    #1
}

\newtcolorbox{solutionbox}[1][]{
    title=Solution / Lösung,
    breakable,
    #1
}

\newtcolorbox{answerbox}[1][]{
    title=Answer / Antwort,
    breakable,
    #1
}

\newtcolorbox{notebox}[1][]{
    title=Notes / Hinweise,
    breakable,
    #1
}

\newtcolorbox{sidecalcbox}[1][]{
    title=Side Calculation / Nebenrechnung,
    breakable,
    #1
}

%%%%%%%%%%%%%%%%%%%%%%%%%%%%%%%%%%%%%%%%%%%%%%%%%%%%%%%%%%%%%%%%%%%
\begin{document}
\paragraph{8. Charged 1D harmonic oscillator in a uniform electric field} We consider a particle of mass $m$ and charge $e$ moving in a (1D) harmonic oscillator potential $V(X) = \frac{1}{2} m \omega^2 X^2$ and additionally in a potential $-e\mathcal{E}X$ generated by a constant electric field $\mathcal{E}$. The Hamiltonian is thus given by \begin{align*} H(\mathcal{E}) = \frac{P^2}{2m} + \frac{1}{2} m \omega^2 X^2 - e\mathcal{E}X. \end{align*} Compute the energy eigenvalues $E_n(\mathcal{E})$ and eigenfunctions $\phi_n(x)$. Discuss the energy shift and the wave functions, and in particular compute the expectation value of the electric dipole moment $\langle D \rangle = e(\phi_n, X \phi_n)$. \textit{Hint:} Complete the potential to a square in $x$ and introduce the new variable $u = x - \frac{e\mathcal{E}}{m\omega^2}$. Show that $u$ and $P$ still satisfy the canonical commutation relations, i.e., they remain a pair of canonically conjugate variables. \begin{calculationbox} Compute Energy spectrum, shifted coordinate, eigenfunctions and dipole moment. $\frac{1}{2}m\omega^2 x^2 - eEx = \frac{1}{2}m\omega^2 \left(x^2 - 2 \frac{eE}{m\omega^2}x\right) = \frac{1}{2} \left( x - \frac{eE}{m\omega^2}\right)^2 - \frac{1}{2}m\omega^2 \left( \frac{eE}{m\omega^2}\right)^2 = \frac{1}{2} \left( x - \frac{eE}{m\omega^2}\right)^2 - \frac{e^2E^2}{2m\omega^2}$ substitution: $u := x - \frac{eE}{m\omega^2}$ $H(E) = \frac{P^2}{2m} + \frac{1}{2} m\omega^2 u^2 - \frac{e^2 E^2}{2 m \omega^2}$ Since $u$ differs by a constant: $[u,P] = [x - \frac{eE}{m\omega^2}, P] = [x,P] - [\frac{eE}{m\omega^2},P] = [x,P] - 0 = i\hbar$. So $u$, $P$ are canonically conjuagte variables. For the normal harmonic scillator the eigenvalues are $E_n^{(0)} \hbar \omega(n + \frac{1}{2})$, $n\in \mathbb{N}_0$ so for our case the eigenvalues are $E_n(E) = \hbar \omega (n + \frac{1}{2} - \frac{e^2E^2}{2m\omega})$ usual eigenfunctionscentered at $x_0=\frac{eE}{m\omega^2}$, so $\phi_n(x;E)=\phi_n\!\left(x-\frac{eE}{m\omega^2}\right)$, where $\phi_n$ is the $n$-th eigenfunction of the standard harmonic oscillator: $\phi_n(u)=\frac{1}{\sqrt{2^n n!}}\left(\frac{m\omega}{\pi\hbar}\right)^{1/4}H_n\left(\sqrt{\frac{m\omega}{\hbar}}\,u\right)e^{-\frac{m\omega}{2\hbar}u^2}$. Therefore $\phi_n(x;E)=\frac{1}{\sqrt{2^n n!}}\left(\frac{m\omega}{\pi\hbar}\right)^{1/4}H_n\left(\sqrt{\frac{m\omega}{\hbar}}\left(x-\frac{eE}{m\omega^2}\right)\right) e^{-\frac{m\omega}{2\hbar}\left(x-\frac{eE}{m\omega^2}\right)^2}$. \end{calculationbox}
\end{document}